\documentclass[11pt]{article}
\usepackage[paperwidth=210mm,paperheight=5000mm,margin=20mm]{geometry}
\usepackage{amsmath,amssymb,xcolor,graphicx}
\usepackage[hypcap=false]{caption}
\usepackage[hidelinks,bookmarksnumbered,bookmarksopen]{hyperref}
\hypersetup{pdftitle={Crystallography — X-ray diffraction and crystal structure},pdfauthor={Lecturer: Dr K. Nakamura},
  pdfcreator={KherveNote}}
\renewcommand\labelitemii{\textbullet}
\renewcommand\labelitemiii{\textbullet}
\renewcommand\labelitemiv{\textbullet}
\renewcommand\labelenumii{\arabic{enumi}.\arabic{enumii}.}
\renewcommand\labelenumiii{\arabic{enumi}.\arabic{enumii}.\arabic{enumiii}.}
\renewcommand\labelenumiv{\arabic{enumi}.\arabic{enumii}.\arabic{enumiii}.\arabic{enumiv}.}
\setlength{\parindent}{0pt}
\setlength{\parskip}{0.6em}
\definecolor{knoteaccent}{HTML}{1A6DD8}
\definecolor{knotemuted}{HTML}{6B6B6B}
\definecolor{knotekey}{HTML}{D96B00}
\newcommand\knotetime[1]{\leavevmode\llap{\color{knotemuted}\scriptsize #1\hspace{1em}}}
\newcommand\knotekey[2][]{\par\noindent#1\colorbox{knotekey!12}{\parbox{\dimexpr\linewidth-2\fboxsep}{%
  \textbf{\color{knotekey}$\star$ Key point.} #2}}\par}
\newcommand\knotequestion[2][]{\par\noindent#1{\color{knoteaccent}\textbf{?}}~\emph{#2}\par}
\newenvironment{knotetranscript}{\par\color{knotemuted}\small}{\par}
\newcommand\knotesummary[1]{\par\noindent\fcolorbox{knoteaccent}{knoteaccent!6}{%
  \parbox{\dimexpr\linewidth-2\fboxsep-2\fboxrule}{\textbf{Summary}\par #1}}\par}
\makeatletter
\newdimen\knote@limit \knote@limit=4950mm
\newbox\knote@page \newbox\knote@chunk
\def\knote@ht#1{\dimexpr\ht#1+\dp#1\relax}
\def\knote@setheight#1{\global\pdfpageheight=#1\relax
  \special{pdf:pagesize width \the\paperwidth\space height \the\pdfpageheight}}
\def\knote@flush{\ifvoid\knote@page\else
  \knote@setheight{\dimexpr\knote@ht\knote@page+40mm+2mm\relax}%
  \box\knote@page\par\clearpage\global\pdfpageheight=\paperheight\fi}
\newenvironment{knotechunk}
  {\par\setbox\knote@chunk=\vbox\bgroup\hsize=\textwidth\linewidth=\textwidth}
  {\par\egroup
   \ifdim\knote@ht\knote@chunk>\knote@limit
     \knote@flush\unvbox\knote@chunk\par\clearpage
   \else\ifvoid\knote@page
     \global\setbox\knote@page=\vbox{\unvbox\knote@chunk}%
   \else\ifdim\dimexpr\knote@ht\knote@page+\knote@ht\knote@chunk\relax>\knote@limit
     \knote@flush\global\setbox\knote@page=\vbox{\unvbox\knote@chunk}%
   \else
     \global\setbox\knote@page=\vbox{\unvbox\knote@page\unvbox\knote@chunk}%
   \fi\fi\fi}
\AtEndDocument{\knote@flush}
\makeatother
\pagestyle{empty}
\begin{document}
\begin{knotechunk}
{\color{knoteaccent}\rule{\linewidth}{1.5pt}}\par
{\LARGE\bfseries Crystallography — X-ray diffraction and crystal structure}\par
{\large Lecturer: Dr K. Nakamura}\hfill{\color{knotemuted}12 October 2026 \textperiodcentered{} Materials building, Room 202}\par
{\color{knoteaccent}\rule{\linewidth}{0.6pt}}\par

\section{Bragg's law}

X-rays reflected from planes $d$ apart interfere constructively when \[n\lambda = 2d\sin\theta\]

\begin{itemize}\setlength{\itemsep}{0pt}\setlength{\parskip}{0pt}
  \item Cu K\ensuremath{\alpha}: $\lambda = 1.5406$ Å (K\ensuremath{\alpha}\textsubscript{1})
  \item a powder pattern plots intensity against $2\theta$
\end{itemize}
\end{knotechunk}

\begin{knotechunk}
\section{Planes and Miller indices}

For a cubic lattice with parameter $a$: \[d_{hkl} = \frac{a}{\sqrt{h^2 + k^2 + l^2}}\]

\subsection{Worked example: copper}

FCC Cu, $a = 3.615$ Å. For (111): $d = 3.615/\sqrt{3} = 2.087$ Å, so $\sin\theta = 1.5406/(2\times2.087) = 0.369$ and $2\theta \approx 43.3^{\circ}$.
\end{knotechunk}

\begin{knotechunk}
\section{Systematic absences}

The structure factor removes some reflections:

\begin{enumerate}\setlength{\itemsep}{0pt}\setlength{\parskip}{0pt}
  \item simple cubic: all $hkl$
  \item body-centred cubic: $h + k + l$ even — (110), (200), (211)…
  \item face-centred cubic: $h, k, l$ all even or all odd — (111), (200), (220), (311)…
\end{enumerate}

\knotekey{The pattern of absences identifies the lattice type before any refinement.}
\end{knotechunk}

\begin{knotechunk}
\section{Peak widths and refinement}

\begin{itemize}\setlength{\itemsep}{0pt}\setlength{\parskip}{0pt}
  \item Scherrer: crystallite size $\tau = \dfrac{K\lambda}{\beta\cos\theta}$, $K \approx 0.9$, $\beta$ the FWHM in radians after removing the instrument broadening
  \item strain also broadens peaks (as $\tan\theta$): Williamson–Hall separates the two
  \item Rietveld refinement fits the whole pattern: lattice parameters, phase fractions, site occupancies
\end{itemize}

\knotequestion{Below what size does Scherrer stop being meaningful, and above which size can it not be used?}
\end{knotechunk}
\end{document}
